% page 585 of the facsimile (image p-584 of the scan)
% block 170.2 x 237.7 mm ; left margin 31.8 mm
\pgc{10.160mm}{0.000mm}{
\l{\cel{56}577}
\l{}
\l{\sou{tercerono}. Persono qua naskis de pelblanko e de mestico. -}
\l{}
\l{\sou{terceto}.\cel{2}Stanco de tri versi. - DEFIR.}
\l{}
\l{\sou{terciana}. (\sou{patol.}) Febro intermitanta, qua rieventas singla-tri-}
\l{\cel{3}dio. - DEFIS.}
\l{}
\l{\sou{terciara}. I. (\sou{geol.)}\cel{2}Dicesas pri un ek la granda faki ek la serio}
\l{\cel{3}geologiala, qua inkluzas la strati eocena, oligocena, miocena e}
\l{\cel{3}pliocena. - II. (\sou{patol.}) Di qua la simptomi aparas pos du}
\l{\cel{3}etapi evidenta di morbo. - DEFIS.}
\l{}
\l{\sou{terebinto}. (\sou{bot.)}\cel{2}Speco di pistachiero. - L.\cel{1}\sou{pistacia terebinthus}.}
\l{\cel{3}- DEFIRS.}
\l{}
\l{\sou{tereno}. Extensajo de tero, de sulo, quan onu egardas kom apta por}
\l{\cel{3}ta o ca uzeso. - DEFIS.}
\l{}
\l{\sou{teriako}. (\sou{medic.}) Elektuario kontre la morduro da la serpenti.-}
\l{\cel{2}DEFIRS.}
\l{}
\l{\sou{teriero}. (\sou{zool.})\cel{2}Mikra hundo quan olim onu uzis por chasar la}
\l{\cel{3}animali qui rezidas en ter-trui. - DEF.}
\l{}
\l{\sou{terino}. Vazo ek ligno, fayenco, porcelano, ceramiko vernisizita,}
\l{\cel{3}ronda, un-peca, sen rebordo nek anso, nek mancho, kun fundo}
\l{\cel{3}plata,uzata en la menaji. - DEFI.}
\l{}
\l{teritorio.\cel{2}Extensajo de lando qua formacas distrikto politikala.}
\l{\cel{3}- DEFIRS.}
\l{}
\l{\sou{termi}. I. (\sou{epoki antiqua}). Establisuro publika por balni varm-aqua.}
\l{\cel{3}- II. (\sou{nun)}. Establisuro adube onu venas por kuracesar per aqui}
\l{\cel{3}medicinala varma. - DEFIS.}
\l{}
\l{\sou{termino}. I. Bloko petra qua indikas la limito di agro, di stadio,}
\l{\cel{3}la frontiero. - II. (\sou{gram.}) En propoziciono : o la subjekto, o}
\l{\cel{3}la verbo, o la atributo. - III. (\sou{logiko}).En silogismo : o la}
\l{\cel{3}majoro, o la minoro, o la mezo. - IV. (matem.)\cel{2}Expresuro di}
\l{\cel{3}quanto en raporto, en proporciono, en progresiono, e c. -}
\l{\cel{3}V. (\sou{mitol.}) La deajo - reprezentita per figuro di viro di qua}
\l{\cel{3}la parto infra finas gaino-forme, skultita en bloko de petro -}
\l{\cel{3}qua sekurigis la egardo di la limiti inter-agra, di la fron-}
\l{\cel{3}tieri, e c. - DEFIS.}
\l{}
\l{\sou{termito}. (\sou{zool.}) Nevroptero qua rodas interne la objekti ek ligno,}
\l{\cel{3}lente e cele. - L.\cel{1}\sou{termes}.\cel{2}- DEFRS.}
\l{}
\l{termodinamiko.\cel{2}Parto di fiziko, qua havas kom studiajo la relati}
\l{\cel{3}qui existas inter la fenomeni mekanikala e kalorifala. - DEFIRS}
\l{}
\l{termoelektro.\cel{2}Ensemblo ek la fenomeni elektrala e kalorala qui}
\l{\cel{3}naskas de la varii od elektrala o kalorala. - DEFIRS}
}
